\begin{abstract}
Se estudia la predicción de empresas ganadoras de licitaciones públicas mediante un modelo de Boltzmann con parámetros comunes a cinco rubros. La energía asignada a cada empresa candidata combina historial de adjudicaciones, similitud textual, menciones del nombre de la empresa y vínculos en una red oficina--empresa. La selección mediante lasso conserva 34 de 44 pistas y alcanza un MRR promedio por rubro de $0{,}282$, frente a $0{,}120$ del modelo de referencia B0; El texto es la principal fuente de mejora, como muestran las ablaciones y los controles aleatorizados. Las correcciones específicas por rubro aportan poco, mientras que la cobertura de la preselección limita el rendimiento.
\end{abstract}
